\title{XQDot4Z: productos punto cuantizados\\en un procesador RISC-V}
\author{\IEEEauthorblockN{Steve Andy Ildefonso Santos}}
\maketitle
\begin{abstract}
Guardar pesos de red neuronal en cuatro bits ahorra memoria, pero el procesador
todavía debe desarmarlos, y la cuantización afín añade un costo fácil de pasar
por alto: un código de peso no es un peso, sino un desplazamiento respecto de un
punto cero que hay que restar primero. Este trabajo pregunta si
esa resta merece construirse dentro de la instrucción de producto punto. Se
añade una instrucción fusionada de cuatro términos a Kuntur, procesador RISC-V
docente de cinco etapas, y se compara contra tres referencias bajo una política
de generación de código congelada: multiplicación por software sin
multiplicador, multiplicación escalar, y un producto punto empacado cuya
corrección se factoriza, calculando la suma de activaciones una vez por grupo y
reutilizándola entre filas. Una campaña registrada de 2280 casos simulados,
todos exitosos en dos simuladores independientes, produce un mapa de condiciones
y no una sola cifra de aceleración. La fusión reduce ciclos frente a ambas,
pero el margen se estrecha al crecer el número de filas, porque la factorizada
reparte su corrección entre más, y desaparece con punto cero cero, donde esa
alternativa la omite. Cuando varía por fila,
la ventaja se paga en tamaño de código: viaja como inmediato, así que cada valor
distinto necesita su propio cuerpo de fila, y la variante fusionada ocupa 516
palabras de instrucción frente a 72. Parte del hueco bruto es control de bucle y
no aritmética, separados midiendo un gemelo estructural. Son ciclos simulados:
no tiempo físico, energía ni área.
\end{abstract}
\begin{IEEEkeywords}
inferencia cuantizada, aritmética de baja precisión, extensiones de instrucciones,
verificación de hardware
\end{IEEEkeywords}
% Keep the abstract first; a top float may start in the next column instead.
\suppressfloats[t]

\section{Introducción}
Una capa de red neuronal es, por debajo, una enorme cantidad de productos punto
entre pesos y activaciones. La cuantización permite guardar ambos lados como
enteros pequeños con un factor de escala~\cite{jacob}, y el ahorro en
almacenamiento es notable: ocho pesos de cuatro bits caben en una palabra de 32
bits. Pero ese es un ahorro de memoria, no de trabajo. El procesador todavía
tiene que desarmar la palabra antes de poder multiplicar algo, y desempacar no
sale gratis.

La cuantización afín añade un segundo costo, más fácil de pasar por alto. Un
código de peso de cuatro bits no es el peso. Es un desplazamiento respecto de un
punto cero compartido por un grupo de pesos, así que lo que la aritmética
necesita de verdad es el código menos ese punto cero. Una resta no es nada.
Cuatro restas en cada grupo de cada producto punto de cada fila de una capa sí
son algo.

Este trabajo hace una pregunta deliberadamente estrecha sobre esa resta: ¿merece
construirse dentro de la instrucción que calcula el producto punto? Esa
instrucción se añade a Kuntur, un procesador RISC-V docente de cinco etapas, y
se mide contra tres puntos de referencia que resuelven la misma aritmética de
maneras distintas. El objetivo no es un acelerador, ni una pila de inferencia
completa. Es un mapa de las condiciones bajo las cuales fusionar la corrección
se paga sola, incluidas aquellas en las que claramente no lo hace.

\begin{figure}[t]
\centering
\begin{tikzpicture}[
 font=\scriptsize,
 box/.style={draw,minimum height=5.2mm,inner xsep=3pt,align=center},
 lbl/.style={font=\scriptsize\bfseries},
 arr/.style={-{Stealth[length=1.4mm]},shorten >=1pt,shorten <=1pt}]

\node[lbl] at (-3.55,0.55) {\ifSpanish Entrada\else Input\fi};
\node[box,text width=17mm] (in) at (-3.55,0)
 {\ifSpanish 8 pesos U4\\4 activaciones S8\else 8 U4 weights\\4 S8 activations\fi};

\foreach \y/\name/\code in {-1.25/b1/B1, -2.15/b2/B2, -3.05/b3/B3, -3.95/d/D}{
  \node[lbl] at (-3.55,\y) {\code};
}

\node[box,text width=20mm] (b1a) at (-1.35,-1.25)
 {\ifSpanish desplazar y sumar\else shift and add\fi};
\node[box,text width=20mm] (b1b) at (1.15,-1.25)
 {\ifSpanish restar 4 veces\else subtract 4 times\fi};

\node[box,text width=20mm] (b2a) at (-1.35,-2.15) {4 $\times$ MUL};
\node[box,text width=20mm] (b2b) at (1.15,-2.15)
 {\ifSpanish restar 4 veces\else subtract 4 times\fi};

\node[box,text width=20mm] (b3a) at (-1.35,-3.05) {XQDot4 $\rightarrow P$};
\node[box,text width=20mm] (b3b) at (1.15,-3.05) {$P - z\,S_a$};

\node[box,text width=20mm,very thick] (da) at (-1.35,-3.95) {XQDot4Zi};
\node[box,text width=20mm,draw=none] (db) at (1.15,-3.95)
 {\ifSpanish (ya corregido)\else (already corrected)\fi};

\node[box,text width=12mm] (out) at (3.35,-2.6) {$d$ (INT32)};

\foreach \a/\b in {b1a/b1b, b2a/b2b, b3a/b3b}{\draw[arr] (\a) -- (\b);}
\draw[arr] (da) -- (db);
\foreach \n in {b1a,b2a,b3a,da}{\draw[arr] (in.east) -| ($(\n.west)+(-3mm,0)$) -- (\n.west);}
\foreach \n in {b1b,b2b,b3b,db}{\draw[arr] (\n.east) -- ++(3mm,0) |- (out.west);}

\node[align=left,text width=26mm,anchor=west] at (2.1,-4.75)
 {\ifSpanish $S_a$ se calcula una vez por grupo y se reutiliza en las $N$ filas.
  \else $S_a$ is computed once per group and reused across the $N$ rows.\fi};
\end{tikzpicture}
\caption{\ifSpanish Trabajo que hace cada variante sobre un mismo grupo de cuatro
pesos. B1--B3 son puntos de comparación; D es el diseño en estudio. Las cuatro
leen los mismos datos empacados. La diferencia no está en el producto sino en
quién paga la corrección del punto cero: B1 y B2 la restan término a término,
B3 la factoriza y la amortiza entre filas, y D la absorbe dentro de la propia
instrucción.
\else Work each variant performs on one group of four weights. B1--B3 are points
of comparison; D is the design under study. All four read the same packed data.
The difference is not in the product but in who pays for the zero-point
correction: B1 and B2 subtract it term by term, B3 factors it and amortizes it
across rows, and D absorbs it inside the instruction itself.\fi}
\label{fig:comparison}
\end{figure}


\section{Antecedentes y trabajo relacionado}
Bajo cuantización afín un código entero \(q\) representa \(r=s(q-z)\), donde
\(s>0\) es una escala y \(z\) un punto cero. Tomando activaciones cuyo punto
cero es cero, el núcleo aritmético en estudio es
\begin{equation}
d=\sum_{i=0}^{3}(q_{w,i}-z_w)q_{a,i}.
\label{eq:dot}
\end{equation}
Los pesos se guardan como U4, enteros sin signo de cuatro bits en \([0,15]\), y
las activaciones como S8, enteros con signo de ocho bits en \([-128,127]\). El
resultado es un subtotal INT32 en \([-7680,7680]\); escalas, sesgo y la
acumulación de otros subtotales quedan fuera de la operación.

Leída al pie de la letra, \eqref{eq:dot} gasta cuatro restas antes de formar
producto alguno. Sin embargo puede reordenarse como
\begin{equation}
d=\sum_i q_{w,i}q_{a,i}-z_w\sum_i q_{a,i}.
\label{eq:factor}
\end{equation}
El primer término es un producto punto corriente sobre los códigos crudos. El
segundo es el punto cero por la suma de las activaciones, y esa suma no depende
de los pesos en absoluto. En un producto matriz--vector las mismas activaciones
se encuentran con todas las filas, así que puede calcularse una vez y
reutilizarse en las \(N\) filas. Esa es justamente la razón por la que una
comparación honesta no puede limitarse a contar las restas de \eqref{eq:dot}:
una implementación competente nunca las paga.

XpulpNN y Dustin son precedentes de extensiones de instrucciones de baja
precisión y precisión mixta sobre RISC-V~\cite{xpulpnn,dustin}. Los productos
punto de pocos bits no son, por tanto, una idea nueva, ni se reclaman como tal
aquí. La pregunta examinada es más estrecha: dado que un producto punto empacado
ya existe, ¿la corrección del punto cero pertenece dentro de la instrucción o
fuera de ella?

\section{Pregunta de investigación y alcance}
\textit{¿Bajo qué condiciones fusionar la corrección del punto cero en un
producto punto U4 por S8 reduce el costo de ejecución frente a multiplicación
escalar y frente a corrección factorizada, considerando la reutilización de
activaciones y el costo de suministrar el punto cero?}

El alcance es un solo núcleo, lote uno, y núcleos pequeños de producto punto y
matriz--vector. Costo significa ciclos de un pipeline in-order simulado,
reportados junto con conteos de instrucciones y tamaño de código. Síntesis,
área, frecuencia y energía pertenecen a una etapa posterior con otra
metodología, y nada de lo medido aquí se extrapola a ellas.

Una distinción recorre todo lo que sigue. Un punto cero puede ser constante en
toda una capa (ZC), variar por fila pero conocerse al generar el código (ZS), o
leerse de memoria en ejecución (ZR). La instrucción propuesta lleva el punto
cero como inmediato de cuatro bits, lo que cubre los dos primeros regímenes y no
el tercero. ZR queda entonces declarado fuera de alcance en vez de estimado,
porque un inmediato no puede sustituir sin costo a un valor leído en ejecución.

\section{Método}
\subsection{Las cuatro variantes}
Cuatro implementaciones producen las mismas salidas a partir de los mismos datos
empacados. Tres son puntos de referencia, nombrados con una B inicial de
\emph{baseline}; la cuarta, D, es el diseño en estudio. La
Tabla~\ref{tab:variants} indica qué las separa y la
Fig.~\ref{fig:comparison} muestra qué hace cada una con un solo grupo de cuatro
pesos.

\begin{table}[t]
\caption{Las cuatro variantes. B marca un punto de referencia y D el diseño en
estudio. Las cuatro leen datos empacados idénticos y comparten una sola política
de generación de código, así que lo que difiere entre ellas es la aritmética, no
el código que la rodea.}
\label{tab:variants}
\centering
\begin{tabular}{llll}
\hline
 & Multiplicador & Producto punto & Corrección del punto cero \\
\hline
B1 & ninguno & descomposición en bits & por término, sin multiplicar \\
B2 & MUL escalar & un MUL por peso & por término \\
B3 & MUL escalar & empacado, cuatro carriles & factorizada, por fila \\
D & no requiere & empacado, cuatro carriles & fusionada en la instrucción \\
\hline
\end{tabular}
\end{table}

B1 no tiene multiplicador alguno y construye cada producto por descomposición en
bits con máscaras acotadas. Es contexto, no una afirmación sobre software
óptimo: responde cuánto cuesta esta operación en un núcleo sin hardware de
multiplicación, y su implementación se congela antes de medir para que no pueda
ajustarse después de ver resultados. B2 es la versión directa con multiplicación
escalar.

B3 es el competidor serio. Corre una unidad empacada de cuatro carriles sobre
los códigos crudos y aplica la corrección en la forma factorizada de
\eqref{eq:factor}, calculando la suma de activaciones una vez por grupo y
reutilizándola entre filas. D usa una unidad de cuatro carriles que resta el
punto cero dentro de cada carril, así que entrega el subtotal ya corregido y
nunca llega a formar una suma de activaciones.

La comparación que responde la pregunta es D contra B3. B1 y B2 acotan el
espacio a su alrededor; no la resuelven.

\subsection{Plataforma y verificación}
El procesador docente original se conserva y todo el desarrollo ocurre en un
clon independiente, de modo que la base no puede desplazarse bajo el
experimento. Cada pieza se verifica por separado antes de combinarse.

La unidad XQDot4Z aislada supera 556\,734 comparaciones por simulador contra una
referencia en Python con enteros no acotados, junto con ocho controles negativos
y cuatro mutaciones inyectadas que las pruebas están obligadas a detectar; el
lint no necesita excepciones. La exhaustividad cubre únicamente tuplas de un
término. Integrarla al pipeline como XQDot4Zi, en su forma con inmediato, añade
25 programas dirigidos en dos simuladores, comparando retiro, writeback, stores
y operandos, y detectando dos mutaciones más.

La MUL escalar opcional que usa B2 es independiente de ese trabajo y devuelve
los 32 bits bajos del producto~\cite{riscvm}; no implementa ni la extensión M
completa ni Zmmul. Supera 86\,704 vectores por simulador y 49 programas en
cuatro configuraciones. La unidad empacada detrás de B3, que omite la resta del
punto cero, supera 67\,910 vectores por simulador; la integración de XQDot4
supera 123 programas en ambos simuladores, cubriendo las ocho configuraciones de
las tres instrucciones opcionales.

Cada hito queda fijado por SHA-256 en un archivo de estado, cada corrida se
repite para mostrar que es determinista, y se inyectan mutaciones a propósito
para demostrar que las pruebas pueden fallar. Dos simuladores independientes,
Icarus Verilog y Verilator, deben coincidir exactamente en cada contador antes
de registrar un caso como exitoso.

\subsection{Cómo se vuelve honesta la comparación}
Una comparación entre una instrucción y su ausencia es fácil de amañar, así que
los grados de libertad se eliminan antes de cualquier medición. Desenrollado,
asignación de registros, planificación y reutilización de la corrección quedan
fijados por una sola política escrita que comparten las cuatro variantes, y esa
política, junto con los generadores y las pruebas que la implementan, se congela
antes de correr la campaña. La política permite explícitamente omitir trabajo
que un punto cero determinado vuelve innecesario, que es como B3 termina
eliminando su corrección por completo cuando el punto cero es cero.

La rejilla, las semillas, los perfiles de punto cero y las reglas de reporte se
registran de antemano, y el registro fija las hipótesis por hash para que no
puedan editarse después. Los regímenes nunca se agregan en una sola cifra, y
todo caso planificado se conserva, fallos incluidos, porque un caso meramente
ausente haría que una campaña incompleta pareciera terminada.

Una sutileza merece mecanismo propio. Como D lleva el punto cero en un
inmediato, un punto cero que varía por fila lo obliga a emitir un cuerpo de fila
por cada fila, mientras que B3 puede conservar un bucle. Parte de cualquier
hueco medido es entonces control de bucle y no corrección. Para separarlos, cada
variante se mide también como gemelo estructural: el mismo cuerpo de bucle
replicado sin la arista de retorno. Restar el gemelo cotiza el control de bucle
con exactitud, y lo que queda es atribuible a la aritmética.

\balance
\section{Resultados y discusión}
La campaña registrada corrió \CampaignCases\ casos, todos exitosos en ambos
simuladores y sin ninguno ausente. Los ciclos son idénticos entre las diez
semillas en las \CampaignCombos\ combinaciones de variante, forma y ajuste, y
entre las dieciséis fases del perfil de punto cero variable, mientras que las
salidas difieren en todas. Eso es una afirmación sobre este núcleo y este
layout de datos, no una prueba de independencia frente a datos arbitrarios, y
explícitamente no se usa para justificar una rejilla más pequeña.

La Fig.~\ref{fig:ratio} y la Tabla~\ref{tab:ratio} contienen el resultado
central. La ventaja de la fusión sobre B3 cae al crecer el número de filas en
todo régimen donde B3 tiene corrección que amortizar, que es la tendencia que el
trabajo esperaba. En \(z{=}0\) no hay ventaja alguna: B3 omite tanto la suma de
activaciones como la corrección, y las dos variantes coinciden ciclo a ciclo. En
\(N{=}1\) un punto cero variable es degenerado, ya que una sola fila no puede
tener un punto cero que varíe, así que esas columnas coinciden necesariamente.

\begin{figure}[!htbp]
\centering
\begin{tikzpicture}
\begin{axis}[
  width=1.02\columnwidth, height=46mm,
  ybar, bar width=3.4pt,
  enlarge x limits=0.28,
  ymin=0.9, ymax=2.05,
  ytick={1.0,1.2,1.4,1.6,1.8,2.0},
  symbolic x coords={1,4,16},
  xtick=data,
  xlabel={\ifSpanish Filas $N$\else Rows $N$\fi},
  ylabel={\ifSpanish Ciclos B3 / ciclos D\else B3 cycles / D cycles\fi},
  tick label style={font=\scriptsize},
  label style={font=\scriptsize},
  legend style={font=\scriptsize,at={(0.5,1.03)},anchor=south,
                legend columns=4,draw=none,column sep=3pt},
  ymajorgrids, grid style={gray!25},
  axis line style={black}, tick style={black},
]
\addplot[fill=white,draw=black] coordinates {(1,1.00) (4,1.00) (16,1.00)};
\addplot[fill=gray!35,draw=black] coordinates {(1,1.86) (4,1.25) (16,1.08)};
\addplot[fill=gray!70,draw=black] coordinates {(1,1.88) (4,1.26) (16,1.08)};
\addplot[fill=black,draw=black] coordinates {(1,1.00) (4,1.46) (16,1.28)};
\legend{$z{=}0$, $z{=}8$, $z{=}15$, ZS}
\draw[dashed,gray] ({rel axis cs:0,0}|-{axis cs:1,1.0})
                -- ({rel axis cs:1,0}|-{axis cs:1,1.0});
\end{axis}
\end{tikzpicture}
\caption{\ifSpanish Ventaja de la fusión por régimen y número de filas. La línea
punteada en 1.0 marca la ausencia de ventaja. Los tres primeros grupos son
puntos cero constantes; ZS los hace variar por fila. La ventaja se estrecha al
crecer $N$ porque B3 amortiza su corrección entre más filas, y en $z{=}0$ no
existe en absoluto porque B3 la omite. En $N{=}1$ un punto cero variable es
degenerado: una sola fila no puede tener un $z$ que varíe. Corrida
\CampaignRun.
\else Advantage of fusion by regime and row count. The dashed line at 1.0 marks
the absence of any advantage. The first three groups are constant zero points;
ZS makes them vary per row. The advantage narrows as $N$ grows because B3
amortizes its correction over more rows, and at $z{=}0$ there is none at all
because B3 omits it. At $N{=}1$ a varying zero point is degenerate: a single row
cannot have a $z$ that varies. Run \CampaignRun.\fi}
\label{fig:ratio}
\end{figure}


Con punto cero variable el hueco bruto mezcla dos efectos, y el gemelo
estructural los separa. De los 164 ciclos que separan a las variantes en
\(N{=}16\), unos 47 pertenecen al control de bucle y 117 a la corrección misma.
El cociente corregido pasa de 1.38 a 1.20, no de 1.46 a 1.28. Reportar solo la
cifra bruta acreditaría a la instrucción un trabajo que hizo la estructura del
bucle.

Frente a MUL escalar, D emplea entre 5.4 y 6.8 veces menos ciclos con
\(N{=}16\), con el extremo bajo en \(z{=}0\), donde D no tiene corrección que
ahorrar; B1, sin multiplicador, gasta entre 26 y 34 veces más que D.

Esa especialización forzada no es gratuita, y el costo aterriza en otra moneda.
Con \(N{=}16\) y punto cero variable, D ocupa 516 palabras de instrucción frente
a 72 de B3, porque su codificación exige dieciséis cuerpos de fila separados. La
ventaja en ciclos y la penalización en tamaño de código son dos caras de la
misma decisión de interfaz: poner el punto cero en un inmediato es lo que vuelve
gratis la corrección en ejecución y lo que hace crecer el código.

\begin{table}[t]
\caption{Ciclos de B3 divididos entre ciclos de D, por régimen y número de
filas. Un valor de 1.00 significa que ambas variantes cuestan exactamente lo
mismo. Corrida \CampaignRun, \CampaignCases\ casos; los regímenes nunca se
agregan.}
\label{tab:ratio}
\centering
\begin{tabular}{crrrr}
\hline
\(N\) & ZC \(z{=}0\) & ZC \(z{=}8\) & ZC \(z{=}15\) & ZS \\
\hline
1 & 1.00 & 1.86 & 1.88 & 1.00 \\
4 & 1.00 & 1.25 & 1.26 & 1.46 \\
16 & 1.00 & 1.08 & 1.08 & 1.28 \\
\hline
\end{tabular}
\end{table}

Leídas en conjunto, estas cifras responden la pregunta con una condición y no
con un veredicto. La fusión ayuda más donde la corrección no puede amortizarse:
pocas filas y punto cero distinto de cero. Ayuda menos, o nada, donde la
alternativa puede repartir la corrección entre muchas filas o saltársela por
completo. Y donde más ayuda en ciclos, con punto cero variable, es donde más
cuesta en código.

\subsection{Amenazas a la validez}
El esqueleto compartido reduce el sesgo de optimización pero no demuestra
optimalidad ni elimina la interacción entre software y arquitectura. La
distribución sintética no representa necesariamente una red real. Con \(z{=}0\)
puede omitirse la corrección separada, lo que no demuestra igualdad de
rendimiento en general. Despachar entre inmediatos sostendría un punto cero
dinámico sin otra codificación, pero aquí no tiene implementación ni costo
medido. Los contadores están validados estructuralmente, contra ocho programas
dirigidos, tres controles negativos, cuatro mutaciones del banco detectadas y
una identidad de contabilidad de ciclos comprobada en cada caso; los ciclos RTL
siguen sin demostrar timing físico. La pregunta de sensibilidad queda contestada
solo para el número de filas: se midió un grupo de 32 elementos, así que cómo se
mueve la comparación con la cantidad de grupos y su tamaño sigue abierto, y el
layout de datos congelado no alcanza hoy esas formas.

\section{Conclusiones}
Fusionar la corrección del punto cero en la instrucción de producto punto reduce
ciclos frente a corrección factorizada y frente a multiplicación escalar, pero
no de manera uniforme. El beneficio se estrecha al crecer el número de filas,
desaparece cuando el punto cero es cero y, cuando el punto cero varía, se cobra
de vuelta en tamaño de código. La respuesta a la pregunta no es un número sino
un mapa de condiciones, y ese mapa es lo que este trabajo aporta.

Quedan fuera el régimen con puntos cero leídos en ejecución, que esta interfaz
no puede expresar, y la evidencia física: área, frecuencia y energía requieren
síntesis y otra metodología. El hardware adicional que exige la variante
fusionada no está contabilizado en ninguna cifra reportada aquí.
